\begin{explanationbox}
	\textbf{\textcolor{CExplanation}{一句话直觉}}
	三角形已知两边及一对角时，解的存在性取决于该角是钝角还是锐角，以及所对边与邻边及其高的大小关系。高 $h = b\sin A$ 是判据的关键。

	\medskip
	\textbf{\textcolor{CExplanation}{核心拆解}}
	\begin{description}[style=nextline, leftmargin=2.8em, labelsep=0.5em, itemsep=0.45em]
		\item[\textbf{要点一（几何高 h）：}]
			记 $h = b\sin A$，它是 $\triangle ABC$ 中 $AB$ 边上的高。当 $A$ 为锐角时，$h$ 是顶点 $C$ 到 $AB$ 的距离；当 $A$ 为钝角时，高在三角形外，但 $h$ 仍表示该距离。
		\item[\textbf{要点二（角分类）：}]
			判据按 $A\geq90^\circ$ 和 $A<90^\circ$ 两类分别讨论。$A$ 为直角或钝角时，$C$ 在 $AB$ 的一侧，仅当 $a>b$ 才有解；$A$ 为锐角时，$C$ 可能在 $AB$ 两侧，解数更丰富。
		\item[\textbf{要点三（大小比较）：}]
			锐角情况下，先比较 $a$ 与 $h$ ，再比较 $a$ 与 $b$。$a<h$ 无解，$a=h$ 唯一解，$h<a<b$ 两解，$a\geq b$ 一解。这是判据的核心结论。
	\end{description}

	\medskip
	\textbf{\textcolor{CExplanation}{几何本质（作图判断）}}
	固定点 $A$，作射线 $AX$ 使得 $\angle BAX = A$，在 $AX$ 上取点 $C$ 使 $AC=b$。以 $C$ 为圆心、$a$ 为半径作圆，该圆与射线 $AB$ 的交点即为点 $B$。交点的个数即为三角形解的个数。$h = b\sin A$ 是点 $C$ 到射线 $AB$ 的距离，即圆与射线的最近距离。

	\medskip
	\textbf{\textcolor{CExplanation}{代数意义（正弦变换）}}
	由正弦定理 $\frac{a}{\sin A}=\frac{b}{\sin B}$ 可得 $\sin B = \frac{b\sin A}{a} = \frac{h}{a}$。方程 $\sin B = \frac{h}{a}$ 在 $(0,\pi)$ 内的解数决定了三角形个数。当 $\frac{h}{a}>1$ 无解；等于 $1$ 一解；小于 $1$ 时需根据 $A$ 的大小和 $a,b$ 关系判断是否两解。

	\medskip
	\textbf{\textcolor{CExplanation}{考点价值（高频应用）}}
	\begin{itemize}[leftmargin=*, itemsep=0.5em, topsep=0.4em]
		\item \textbf{考法一（直接判解数）：} 给出 $a,b,A$，直接利用判据确定三角形解的个数。常见于选择题或填空题。
		\item \textbf{考法二（多解讨论）：} 在解三角形大题中，已知 $a,b,A$ 求角 $B$，由 $a<b$ 且 $A$ 锐角时可能出现两个解，需要检验大边对大角判定取舍。
	\end{itemize}

	\medskip
	\textbf{\textcolor{CExplanation}{顿悟点}}
	\textbf{$h = b\sin A$ 是判据的桥梁。锐角情况下，$a$ 与 $h$ 的比较决定基础解数（无解、唯一解还是可能两解），$a$ 与 $b$ 的比较再决定是否两解。记住先比 $h$ 再比 $b$。}

	\medskip
	\textbf{\textcolor{CExplanation}{使用场景}}
	\begin{itemize}[leftmargin=*, itemsep=0.5em, topsep=0.4em]
		\item \textbf{场景一（SSA 求角）：} 已知 $a,b,A$，求角 $B$ 或边 $c$，先用判据确定解的个数，再计算。
		\item \textbf{场景二（解的验证）：} 解得 $B$ 后，利用补角是否满足三角形内角和进行取舍，本质上与判据一致。
	\end{itemize}

\end{explanationbox}
